% Første linje er dokumentklassen. Brug *altid* memoir!
\documentclass[a4paper,oneside,article]{memoir}

% Herunder indlæses forskellige pakker
\usepackage[utf8]{inputenc}  % Korrekt håndtering af æ, ø og å
\usepackage[T1]{fontenc}  % Korrekt håndtering af æ, ø og å
\usepackage{microtype}  % Typografisk magi! Giver bl.a. pænere orddeling
\usepackage{graphicx}  % Gør det muligt at indsætte billeder
\usepackage{amsmath}  % Giver adgang til uundværlige matematikting
\usepackage{mathtools}  % Retter ting i ams og giver ekstra muligheder
\usepackage[danish]{babel}  % Danske betegnelser og orddeling
\usepackage{lmodern}
\usepackage{alphabeta}
\usepackage{float} %Mere kontrol over billede-placering
\renewcommand{\danishhyphenmins}{22}  % Bedre dansk orddeling
\usepackage[exponent-product=\ensuremath{\cdot}]{siunitx} % Sientific notation for forskellige ting, med en regel om at den skal bruge cdot istedet for times. 
\usepackage{derivative} % Til at skrive differentialligninger 
\usepackage{xcolor} % For at få flere tekst farver 
\usepackage[colorlinks=true]{hyperref} % For at få refrences man kan trykke på, slet [] hvis man vil have kasser i stedet for farvet tekst.
\usepackage{pdfpages} % Til at indsætte andre pdf'er
\usepackage{subfig} % For at have to billeder ved siden af hinanden
\usepackage{unicode-math}
\usepackage{bm}




% Pænere toc indents 
\usepackage{tocloft}
\cftsetindents{section}{1em}{2em}
\cftsetindents{subsection}{4em}{0em}


% Indholdet i dit dokument starter her!
\begin{document}
\author{Mikkel Moth Billing \\ 202307989}
\title{Atom- og Molekylefysik Formelsamling}
\date{\today}
\maketitle   % Indsætter overskriften
% Table of Contents indstillinger
\setcounter{tocdepth}{2}
\setcounter{secnumdepth}{1}
\hypersetup{linkcolor=black} % Laver tableofcontents sort.  
\tableofcontents %Indsætter ToC
\hypersetup{linkcolor=red} % og så alle andre hyperrefs røde. 

%Custom Commands
%Til opgavebeskrivelser
\newcommand{\opg}[1]{\textcolor{gray}{\emph{#1}}\plainbreak{1}} 

\newcommand{\ti}[1]{\cdot 10^{#1}} % Let: gange 10^x
\newcommand{\e}[1]{\emph{e}^{#1}} % Pænere e^x
\newcommand{\vm}[2]{\begin{pmatrix} #1 \\ #2 \end{pmatrix}}
\newcommand{\mat}[1]{\begin{bmatrix} #1 \end{bmatrix}}

% Lette paranteser 
\newcommand{\br}[2][1]{%
\ifcase#1
\or \left( #2 \right) %1 eller ingenting
\or \left[ #2 \right] %2
\or \left\langle #2 \right\rangle %3
\or \left\{ #2 \right\} %4
\or \left| #2 \right| %5
\fi
}
%Til at sætter bileder ind i store dokumenter
\newcommand{\bil}[1]{
    \begin{figure}[H]
        \centering
        \includegraphics[width=0.75\linewidth]{Billeder/#1.png}
        \label{fig:#1}
    \end{figure}
    }
    
%Til at sætte pdf'er ind i store dokumenter (kun pdf'er på mere end 1 sider)
\newcommand{\pdf}[4]{
    \includepdf[pages=#3,scale=.85,pagecommand={\subsection{#2 - \ref{ch:opg}}  \label{asc:#2}}]{pdf/#1.pdf}
    \includepdf[pages=\the\numexpr 1 + #3\relax-#4, scale=.85]{pdf/#1.pdf}
}

\newcommand{\h}{\hbar}
\newcommand{\dd}{\partial}
\newcommand{\eq}[1]{\label{eq#1}\tag{#1}}
\renewcommand{\a}{\hat{a}}
\newcommand{\n}{\plainbreak{1}}
\newcommand{\ep}{\epsilon}
\renewcommand{\r}{\text{\usefont{U}{dutchcal}{m}{n}r}}


\newcommand{\hv}[2][]{%
    \if\relax\detokenize{#1}\relax
        |#2\rangle%
    \else
        \if\relax\detokenize{#2}\relax
            \langle #1|%
        \else
            \langle #1|#2\rangle%
        \fi
    \fi
}
 
 
% Lav ny pagestyle (kun odd fordi enkeltsidet)
\makepagestyle{Title}
\pagestyle{Title}
%\makeoddhead{Title}{\footnotesize\color{gray}{\theauthor}}{}{\textcolor{gray}{\thedate}}
\makeoddfoot{Title}{}{\footnotesize\color{gray}{\thepage{} / \thelastpage}}{}
% Også fod på siden med \maketitle
\makeoddfoot{plain}{}{\footnotesize\color{gray}{\thepage{} / \thelastpage}}{}
\newpage
 
%Skriv
\chapter{Sandwichformlen}
In quantum mechanics, dipole transitions involve matrix elements of the form:
\begin{equation}
    \langle \psi_f | r_i | \psi_i \rangle,
\end{equation}
where $r_i = x, y,$ or $z$ represents the position operator components in spherical coordinates.

The goal is to determine which operator ($x, y,$ or $z$) to use based on selection rules and symmetry.

Step 1: Check the Selection Rules
Dipole transitions must satisfy the following:
\begin{enumerate}
    \item \textbf{Change in orbital angular momentum:} 
    \begin{equation}
        \Delta l = \pm1.
    \end{equation}
    \item \textbf{Change in magnetic quantum number $m$:}
    \begin{equation}
        \Delta m = 0, \pm1.
    \end{equation}
    \item Different operators couple different values of $m$:
    \begin{itemize}
        \item $x$ and $y$ couple $m \to m \pm 1$.
        \item $z$ couples $m \to m$ (i.e., $\Delta m = 0$).
    \end{itemize}
\end{enumerate}

Step 2: Transformation of $x, y, z$ in Terms of Spherical Harmonics.
In spherical coordinates, the position operators transform as:
\begin{align}
    x &= r \sin\theta \cos\phi \quad \text{(behaves like $Y_1^{+1} + Y_1^{-1}$)} \\
    y &= r \sin\theta \sin\phi \quad \text{(behaves like $i(Y_1^{+1} - Y_1^{-1})$)} \\
    z &= r \cos\theta \quad \text{(behaves like $Y_1^0$)}.
\end{align}

This means:
\begin{itemize}
    \item $z$ can only couple states with $\Delta m = 0$.
    \item $x$ and $y$ couple states with $\Delta m = \pm1$.
    \item $y$ introduces an imaginary factor $i$, while $x$ does not.
\end{itemize}

Step 3: Match $\Delta m$ to the Right Operator. 
The choice of operator depends on the magnetic quantum number change:

\begin{table}[h]
    \centering
    \begin{tabular}{|c|c|c|}
        \hline
        \textbf{Initial $m$} & \textbf{Final $m$} & \textbf{Use Operator} \\
        \hline
        0 & 0 & $z$ \\
        0 & $\pm1$ & $x$ or $y$ \\
        $\pm1$ & 0 & $x$ or $y$ \\
        $\pm1$ & $\pm1$ & Forbidden \\
        $\pm1$ & $\mp1$ & $x$ or $y$ \\
        Any & $\pm2$ & Forbidden \\
        \hline
    \end{tabular}
    \caption{Operator Selection Based on $\Delta m$}
\end{table}

Step 4: Check for Integral Symmetry. 
Even if the selection rules allow a transition, the integral may still vanish due to symmetry:
\begin{itemize}
    \item If the integral contains an \textbf{odd function over a symmetric range}, it is \textbf{zero}.
    \item If the integral contains an \textbf{orthogonal spherical harmonic combination}, it is \textbf{zero}.
    \item If the integral using $x$ is zero, try $y$ instead. The extra \textbf{imaginary factor} in $y$ can sometimes fix cancellations.
\end{itemize}

Example Applications
Example 1: $\langle \psi_{100} | x | \psi_{211} \rangle$
\begin{itemize}
    \item $\psi_{100}$ has $m = 0$.
    \item $\psi_{211}$ has $m = +1$.
    \item $\Delta m = +1$ $\Rightarrow$ Use $x$ or $y$.
    \item The $x$-integral vanishes, but $y$ introduces an $i$ factor and gives a nonzero result.
\end{itemize}
\textbf{Final choice:} Use $y$ instead of $x$.

Example 2: $\langle \psi_{100} | z | \psi_{210} \rangle$
\begin{itemize}
    \item $\psi_{100}$ has $m = 0$.
    \item $\psi_{210}$ has $m = 0$.
    \item $\Delta m = 0$ $\Rightarrow$ Use $z$.
    \item The integral is nonzero.
\end{itemize}
\textbf{Final choice:} Use $z$.

Example 3: $\langle \psi_{100} | y | \psi_{21-1} \rangle$
\begin{itemize}
    \item $\psi_{100}$ has $m = 0$.
    \item $\psi_{21-1}$ has $m = -1$.
    \item $\Delta m = -1$ $\Rightarrow$ Use $x$ or $y$.
    \item The $x$-integral vanishes, but $y$ introduces an imaginary phase.
\end{itemize}
\textbf{Final choice:} Use $y$.

Summary
\begin{enumerate}
    \item \textbf{Check selection rules:} Ensure $\Delta l = \pm1$ and $\Delta m = 0, \pm1$.
    \item \textbf{Match $\Delta m$ to the correct operator:}
    \begin{itemize}
        \item $\Delta m = 0 \Rightarrow$ Use $z$.
        \item $\Delta m = \pm1 \Rightarrow$ Use $x$ or $y$.
    \end{itemize}
    \item \textbf{Check spherical harmonics for integral symmetry:} If $x$-integral is zero, try $y$.
    \item \textbf{Use the decision table to quickly determine the valid operator.}
\end{enumerate}

\end{document}